% !PRINTED none
% Abschnitt-opening page: no folio number visible in the running head on this page
% (chapter-opening layout). Ninth Chapter / Discriminant / § 84 begins here.

\begin{center}
\gesperrt{Ninth Chapter.}

\textbf{\Large Discriminant.}

\rule{0.3\linewidth}{0.4pt}
\end{center}

\begin{center}
\textbf{§ 84. \gesperrt{Definition of the Discriminants.}}
\end{center}

In the theory of quadratic fields there occur as discriminants certain rational integers (positive or negative) which, if they are even, are divisible by 4, and which, if they are odd, leave the residue 1 on division by 4.

\gesperrt{Kronecker} has called such numbers \glqq \gesperrt{numbers of discriminant form}\grqq{}. We shall here call them briefly discriminants, and shall not take offence at the fact that this word has still various other meanings in algebra. We therefore define:

1. \gesperrt{An integer, positive or negative and different from zero, which satisfies one of the two congruences}
\weq{1}{D\equiv 0,\quad D\equiv 1 \quad (\mathrm{mod}\ 4)}
\gesperrt{is called a discriminant.}

By this definition every square number is a discriminant. But since these occupy a certain exceptional position, we shall call them \gesperrt{improper discriminants whenever the distinction is necessary.}

2. \gesperrt{The product of two, and consequently of arbitrarily many, discriminants is again a discriminant.}

3. \gesperrt{A discriminant which contains (apart from 1) no square divisor after the separation of which a discriminant remains is called a fundamental discriminant.}

In what follows, fundamental discriminants shall be denoted by $\Delta$, to distinguish them from others. If $D$ is not a fundamental discriminant, then there is one and only one square number $Q^2$ and one fundamental discriminant $\Delta$ such that
\weq{2}{D = \Delta Q^2}
holds; $\Delta$ is then called the \gesperrt{fundamental discriminant associated with} $D$.
